\IfFileExists{stacks-project.cls}{%
\documentclass{stacks-project}
}{%
\documentclass{amsart}
}

% For dealing with references we use the comment environment
\usepackage{verbatim}
\newenvironment{reference}{\comment}{\endcomment}
%\newenvironment{reference}{}{}
\newenvironment{slogan}{\comment}{\endcomment}
\newenvironment{history}{\comment}{\endcomment}

% For commutative diagrams we use Xy-pic
\usepackage[all]{xy}

% We use 2cell for 2-commutative diagrams.
\xyoption{2cell}
\UseAllTwocells

% We use multicol for the list of chapters between chapters
\usepackage{multicol}

% This is generally recommended for better output
\usepackage{lmodern}
\usepackage[T1]{fontenc}

% For cross-file-references
\usepackage{xr-hyper}

% Package for hypertext links:
\usepackage{hyperref}

% For any local file, say "hello.tex" you want to link to please
% use \externaldocument[hello-]{hello}
\externaldocument[introduction-]{introduction}
\externaldocument[conventions-]{conventions}
\externaldocument[sets-]{sets}
\externaldocument[categories-]{categories}
\externaldocument[topology-]{topology}
\externaldocument[sheaves-]{sheaves}
\externaldocument[sites-]{sites}
\externaldocument[stacks-]{stacks}
\externaldocument[fields-]{fields}
\externaldocument[algebra-]{algebra}
\externaldocument[brauer-]{brauer}
\externaldocument[homology-]{homology}
\externaldocument[derived-]{derived}
\externaldocument[simplicial-]{simplicial}
\externaldocument[more-algebra-]{more-algebra}
\externaldocument[smoothing-]{smoothing}
\externaldocument[modules-]{modules}
\externaldocument[sites-modules-]{sites-modules}
\externaldocument[injectives-]{injectives}
\externaldocument[cohomology-]{cohomology}
\externaldocument[sites-cohomology-]{sites-cohomology}
\externaldocument[dga-]{dga}
\externaldocument[dpa-]{dpa}
\externaldocument[sdga-]{sdga}
\externaldocument[hypercovering-]{hypercovering}
\externaldocument[schemes-]{schemes}
\externaldocument[constructions-]{constructions}
\externaldocument[properties-]{properties}
\externaldocument[morphisms-]{morphisms}
\externaldocument[coherent-]{coherent}
\externaldocument[divisors-]{divisors}
\externaldocument[limits-]{limits}
\externaldocument[varieties-]{varieties}
\externaldocument[topologies-]{topologies}
\externaldocument[descent-]{descent}
\externaldocument[perfect-]{perfect}
\externaldocument[more-morphisms-]{more-morphisms}
\externaldocument[flat-]{flat}
\externaldocument[groupoids-]{groupoids}
\externaldocument[more-groupoids-]{more-groupoids}
\externaldocument[etale-]{etale}
\externaldocument[chow-]{chow}
\externaldocument[intersection-]{intersection}
\externaldocument[pic-]{pic}
\externaldocument[weil-]{weil}
\externaldocument[adequate-]{adequate}
\externaldocument[dualizing-]{dualizing}
\externaldocument[duality-]{duality}
\externaldocument[discriminant-]{discriminant}
\externaldocument[derham-]{derham}
\externaldocument[local-cohomology-]{local-cohomology}
\externaldocument[algebraization-]{algebraization}
\externaldocument[curves-]{curves}
\externaldocument[resolve-]{resolve}
\externaldocument[models-]{models}
\externaldocument[functors-]{functors}
\externaldocument[equiv-]{equiv}
\externaldocument[pione-]{pione}
\externaldocument[etale-cohomology-]{etale-cohomology}
\externaldocument[proetale-]{proetale}
\externaldocument[relative-cycles-]{relative-cycles}
\externaldocument[more-etale-]{more-etale}
\externaldocument[trace-]{trace}
\externaldocument[crystalline-]{crystalline}
\externaldocument[spaces-]{spaces}
\externaldocument[spaces-properties-]{spaces-properties}
\externaldocument[spaces-morphisms-]{spaces-morphisms}
\externaldocument[decent-spaces-]{decent-spaces}
\externaldocument[spaces-cohomology-]{spaces-cohomology}
\externaldocument[spaces-limits-]{spaces-limits}
\externaldocument[spaces-divisors-]{spaces-divisors}
\externaldocument[spaces-over-fields-]{spaces-over-fields}
\externaldocument[spaces-topologies-]{spaces-topologies}
\externaldocument[spaces-descent-]{spaces-descent}
\externaldocument[spaces-perfect-]{spaces-perfect}
\externaldocument[spaces-more-morphisms-]{spaces-more-morphisms}
\externaldocument[spaces-flat-]{spaces-flat}
\externaldocument[spaces-groupoids-]{spaces-groupoids}
\externaldocument[spaces-more-groupoids-]{spaces-more-groupoids}
\externaldocument[bootstrap-]{bootstrap}
\externaldocument[spaces-pushouts-]{spaces-pushouts}
\externaldocument[spaces-chow-]{spaces-chow}
\externaldocument[groupoids-quotients-]{groupoids-quotients}
\externaldocument[spaces-more-cohomology-]{spaces-more-cohomology}
\externaldocument[spaces-simplicial-]{spaces-simplicial}
\externaldocument[spaces-duality-]{spaces-duality}
\externaldocument[formal-spaces-]{formal-spaces}
\externaldocument[restricted-]{restricted}
\externaldocument[spaces-resolve-]{spaces-resolve}
\externaldocument[formal-defos-]{formal-defos}
\externaldocument[defos-]{defos}
\externaldocument[cotangent-]{cotangent}
\externaldocument[examples-defos-]{examples-defos}
\externaldocument[algebraic-]{algebraic}
\externaldocument[examples-stacks-]{examples-stacks}
\externaldocument[stacks-sheaves-]{stacks-sheaves}
\externaldocument[criteria-]{criteria}
\externaldocument[artin-]{artin}
\externaldocument[quot-]{quot}
\externaldocument[stacks-properties-]{stacks-properties}
\externaldocument[stacks-morphisms-]{stacks-morphisms}
\externaldocument[stacks-limits-]{stacks-limits}
\externaldocument[stacks-cohomology-]{stacks-cohomology}
\externaldocument[stacks-perfect-]{stacks-perfect}
\externaldocument[stacks-introduction-]{stacks-introduction}
\externaldocument[stacks-more-morphisms-]{stacks-more-morphisms}
\externaldocument[stacks-geometry-]{stacks-geometry}
\externaldocument[moduli-]{moduli}
\externaldocument[moduli-curves-]{moduli-curves}
\externaldocument[examples-]{examples}
\externaldocument[exercises-]{exercises}
\externaldocument[guide-]{guide}
\externaldocument[desirables-]{desirables}
\externaldocument[coding-]{coding}
\externaldocument[obsolete-]{obsolete}
\externaldocument[fdl-]{fdl}
\externaldocument[index-]{index}

% Theorem environments.
%
\theoremstyle{plain}
\newtheorem{theorem}[subsection]{Theorem}
\newtheorem{proposition}[subsection]{Proposition}
\newtheorem{lemma}[subsection]{Lemma}

\theoremstyle{definition}
\newtheorem{definition}[subsection]{Definition}
\newtheorem{example}[subsection]{Example}
\newtheorem{exercise}[subsection]{Exercise}
\newtheorem{situation}[subsection]{Situation}

\theoremstyle{remark}
\newtheorem{remark}[subsection]{Remark}
\newtheorem{remarks}[subsection]{Remarks}

\numberwithin{equation}{subsection}

% Macros
%
\def\lim{\mathop{\mathrm{lim}}\nolimits}
\def\colim{\mathop{\mathrm{colim}}\nolimits}
\def\Spec{\mathop{\mathrm{Spec}}}
\def\Hom{\mathop{\mathrm{Hom}}\nolimits}
\def\Ext{\mathop{\mathrm{Ext}}\nolimits}
\def\SheafHom{\mathop{\mathcal{H}\!\mathit{om}}\nolimits}
\def\SheafExt{\mathop{\mathcal{E}\!\mathit{xt}}\nolimits}
\def\Sch{\mathit{Sch}}
\def\Mor{\mathop{\mathrm{Mor}}\nolimits}
\def\Ob{\mathop{\mathrm{Ob}}\nolimits}
\def\Sh{\mathop{\mathit{Sh}}\nolimits}
\def\NL{\mathop{N\!L}\nolimits}
\def\CH{\mathop{\mathrm{CH}}\nolimits}
\def\proetale{{pro\text{-}\acute{e}tale}}
\def\etale{{\acute{e}tale}}
\def\QCoh{\mathit{QCoh}}
\def\Ker{\mathop{\mathrm{Ker}}}
\def\Im{\mathop{\mathrm{Im}}}
\def\Coker{\mathop{\mathrm{Coker}}}
\def\Coim{\mathop{\mathrm{Coim}}}

% Boxtimes
%
\DeclareMathSymbol{\boxtimes}{\mathbin}{AMSa}{"02}

%
% Macros for moduli stacks/spaces
%
\def\QCohstack{\mathcal{QC}\!\mathit{oh}}
\def\Cohstack{\mathcal{C}\!\mathit{oh}}
\def\Spacesstack{\mathcal{S}\!\mathit{paces}}
\def\Quotfunctor{\mathrm{Quot}}
\def\Hilbfunctor{\mathrm{Hilb}}
\def\Curvesstack{\mathcal{C}\!\mathit{urves}}
\def\Polarizedstack{\mathcal{P}\!\mathit{olarized}}
\def\Complexesstack{\mathcal{C}\!\mathit{omplexes}}
% \Pic is the operator that assigns to X its picard group, usage \Pic(X)
% \Picardstack_{X/B} denotes the Picard stack of X over B
% \Picardfunctor_{X/B} denotes the Picard functor of X over B
\def\Pic{\mathop{\mathrm{Pic}}\nolimits}
\def\Picardstack{\mathcal{P}\!\mathit{ic}}
\def\Picardfunctor{\mathrm{Pic}}
\def\Deformationcategory{\mathcal{D}\!\mathit{ef}}
\title{Units of a monoidal category correspond to pairs making the tensoring functors equivalences}
\author{The Stacks Project Authors}
\date{stacks-project:categories.tex:local-snapshot-20261012; GNU FDL 1.2 or later}
\begin{document}
\maketitle
\noindent\textit{Editorial note.} Counted span is the author's own proof environment, source lines 9240-9318, following the statement at 9230-9238.\par\medskip
\noindent\textit{Editorial note.} Source typo at line 9241: 'isomorpism' for 'isomorphism' (reproduced as-is; no mathematical content affected).\par\medskip
\noindent\textit{Allowed standard prerequisites (editorial).}\begin{itemize}
\item Monoidal category and unit coherence equation. Definition of monoidal category and unit in this chapter; the coherence equation displayed immediately before the lemma (around lines 9200-9226). Fixes what a unit (1, l, r) must satisfy; the displayed equation is the coherence the constructed l and r must reproduce.
\item Equivalences are fully faithful and essentially surjective. Standard framework fact, used without citation at proof lines 9244-9248. L(l\_X) = 1 x id\_X pins l\_X uniquely because L is fully faithful; essential surjectivity gives the reduction to X = Z x 1 and Y = 1 x W.
\end{itemize}\medskip
\begin{lemma}
\label{lemma-monoidal-unit}
Let $(\mathcal{C}, \otimes, \phi)$ be as above. There is a 1-to-1
correspondence between units $(\mathbf{1}, l, r)$ in $\mathcal{C}$ and pairs
$(\mathbf{1}, 1)$ where $\mathbf{1}$ is an object of $\mathcal{C}$ and
$1 : \mathbf{1} \otimes \mathbf{1} \to \mathbf{1}$
is an isomorphism such that the functors $L : X \mapsto \mathbf{1} \otimes X$
and $R : X \mapsto X \otimes \mathbf{1}$ are equivalences.
\end{lemma}

\begin{proof}
Given a unit $(\mathbf{1}, l, r)$ we get an isomorpism
$r : \mathbf{1} \otimes \mathbf{1} \to \mathbf{1}$ and $L$ and $R$ are
equivalences as they are isomorphic to the identity functor.
Conversely, suppose given $(\mathbf{1}, 1)$ such that
$L$ and $R$ are equivalences. We obtain functorial isomorphisms
$l_X : \mathbf{1} \otimes X \to X$ and $r_X : X \otimes \mathbf{1} \to X$
characterized by $L(l_X) = 1 \otimes \text{id}_X$ and
$R(r_X) = \text{id}_X \otimes 1$. Then we have to show that the two
arrows $\text{id}_X \otimes l_Y$ and $r_X \otimes \text{id}_Y$ from
$X \otimes \mathbf{1} \otimes Y$ to $X \otimes Y$ are the same
for all $X$ and $Y$. This property only depends on the isomorphism
classes of $X$ and $Y$.
Since $R$ and $L$ are equivalences, it suffices to
do this for $X = Z \otimes \mathbf{1}$ and $Y = \mathbf{1} \otimes W$
for some objects $Z$ and $W$.
In other words, we have to show that
$$
\text{id}_Z \otimes \text{id}_\mathbf{1} \otimes l_{\mathbf{1} \otimes W}
=
r_{Z \otimes \mathbf{1}} \otimes \text{id}_\mathbf{1} \otimes \text{id}_Y
$$
By construction these maps are equal to
$\text{id}_Z \otimes 1 \otimes \text{id}_\mathbf{1} \otimes \text{id}_W$ and
$\text{id}_Z \otimes \text{id}_\mathbf{1} \otimes 1 \otimes \text{id}_W$.
Thus it suffices to show that
$1 \otimes \text{id}_\mathbf{1} = \text{id}_\mathbf{1} \otimes 1$.

\medskip\noindent
We have
$\text{id}_\mathbf{1} \otimes 1 = r_\mathbf{1} \otimes \text{id}_\mathbf{1}$
and
$l_\mathbf{1} \otimes \text{id}_\mathbf{1} = \text{id}_\mathbf{1} \otimes 1$.
We may write $r_\mathbf{1} = a \circ 1$ and $l_\mathbf{1} = b \circ 1$
for some $a, b$ automorphisms of $\mathbf{1}$. Thus we have
$\text{id}_\mathbf{1} \otimes 1 =
(a \otimes \text{id}_\mathbf{1}) \circ (1 \otimes \text{id}_\mathbf{1})$
and
$1 \otimes \text{id}_\mathbf{1} =
(\text{id}_\mathbf{1} \otimes b) \circ (\text{id}_\mathbf{1} \otimes 1)$.
Then we can write
\begin{align*}
(1 \otimes \text{id}_\mathbf{1}) \circ
(\text{id}_\mathbf{1} \otimes 1 \otimes \text{id}_\mathbf{1})
& =
(1 \otimes \text{id}_\mathbf{1}) \circ
(\text{id}_\mathbf{1} \otimes \text{id}_\mathbf{1} \otimes b)
\circ
(\text{id}_\mathbf{1} \otimes \text{id}_\mathbf{1} \otimes 1) \\
& =
(\text{id}_\mathbf{1} \otimes b) \circ
(1 \otimes \text{id}_\mathbf{1}) \circ
(\text{id}_\mathbf{1} \otimes \text{id}_\mathbf{1} \otimes 1) \\
& =
(\text{id}_\mathbf{1} \otimes b) \circ
(1 \otimes 1)
\end{align*}
and we also have
\begin{align*}
(1 \otimes \text{id}_\mathbf{1}) \circ
(\text{id}_\mathbf{1} \otimes 1 \otimes \text{id}_\mathbf{1})
& =
(\text{id}_\mathbf{1} \otimes b) \circ
(\text{id}_\mathbf{1} \otimes 1) \circ
(a \otimes \text{id}_\mathbf{1} \otimes \text{id}_\mathbf{1}) \circ
(1 \otimes \text{id}_\mathbf{1} \otimes \text{id}_\mathbf{1}) \\
& =
(\text{id}_\mathbf{1} \otimes b) \circ
(a \otimes \text{id}_\mathbf{1}) \circ
(1 \otimes 1)
\end{align*}
This proves that $a \otimes \text{id}_\mathbf{1}$ is the identity and
hence $a$ is the identity as desired.

\medskip\noindent
To finish the proof, we note that the rules above determine
inverse equivalences of categories between the category of units
(suitably defined) and the category of pairs $(\mathbf{1}, 1)$.
\end{proof}

\end{document}
